\documentclass[12pt,a4paper]{article}

\usepackage[spanish,es-nodecimaldot]{babel}
\usepackage[utf8]{inputenc}
\usepackage[T1]{fontenc}
\usepackage{lmodern}
\usepackage{geometry}
\usepackage{microtype}
\usepackage{setspace}
\usepackage{parskip}
\usepackage{enumitem}
\usepackage{xcolor}
\usepackage{listings}
\usepackage{hyperref}
\usepackage{amsmath}
\usepackage{amssymb}
\usepackage{array}
\usepackage{booktabs}
\usepackage{longtable}

\geometry{margin=2.5cm}
\onehalfspacing
\setlist[itemize]{leftmargin=*,itemsep=0.2em}
\setlist[enumerate]{leftmargin=*,itemsep=0.2em}

\hypersetup{
  colorlinks=true,
  linkcolor=black,
  urlcolor=blue,
  pdftitle={Resolucion del Practico 6 - Testing de Software 2024},
  pdfauthor={}
}

\lstdefinestyle{codestyle}{
  basicstyle=\ttfamily\small,
  frame=single,
  breaklines=true,
  showstringspaces=false,
  keywordstyle=\color{blue!70!black}\bfseries,
  commentstyle=\color{gray!70!black}\itshape,
  stringstyle=\color{orange!70!black}
}

\title{\textbf{Resolucion del Practico 6}\\Testing de Software 2024}
\author{}
\date{}

\begin{document}

\maketitle
\tableofcontents
\newpage

\section*{Introduccion}
\addcontentsline{toc}{section}{Introduccion}

Este documento contiene la resolucion del \textit{Practico 6} de Testing de Software, centrado
en dos tecnicas modernas de generacion y evaluacion de tests:

\begin{itemize}
  \item \textbf{Testing de mutacion} con la herramienta \texttt{Pitest}. En los Ejercicios 1 y 2
  se generan mutantes sobre el codigo bajo prueba (\texttt{Palindrome.capicua} y
  \texttt{TriTyp.triang}), se construye una suite JUnit que pretende matarlos a todos y se
  analiza el reporte de PIT (score, sobrevivientes, equivalentes).
  \item \textbf{Fuzzing} en el Ejercicio 3. Se completan los componentes \texttt{Mutator},
  \texttt{MutationFuzzer} y \texttt{RandomFuzzer} y se cierra un test parametrizado que ejecuta
  el comando \texttt{bc} de Linux usando entradas generadas por mutacion.
\end{itemize}

\textbf{Convenciones.}

\begin{itemize}
  \item \texttt{KILLED}: mutante muerto por al menos un test.
  \item \texttt{SURVIVED}: mutante sobreviviente (no es matado por ningun test).
  \item \emph{Mutation Score} bruto: $\frac{\#KILLED}{\#GENERATED}$.
  \item \emph{Mutante equivalente}: mutante cuyo comportamiento observable es identico al del
  programa original; no puede matarse y por eso no debe contar contra el score.
\end{itemize}

\textbf{Como se ejecuta PIT.} Desde \texttt{tp6/assignment-6-rodeghiero}:

\begin{lstlisting}[style=codestyle]
JAVA_HOME=$(/usr/libexec/java_home -v 17) mvn -Dmaven.repo.local=.m2 \
    -Djacoco.skip=true test org.pitest:pitest-maven:mutationCoverage
\end{lstlisting}

La flag \texttt{-Djacoco.skip=true} desactiva el agente de JaCoCo (incompatible con la JVM 17
con la version del plugin del template); ademas, PIT instrumenta el bytecode por su cuenta, asi
que no se pierde nada al saltearlo.

\section{Ejercicio 1: \texttt{Palindrome.capicua}}

\subsection{Codigo bajo prueba}

\begin{lstlisting}[style=codestyle,language=Java]
public static boolean capicua(char[] list) {
    int index = 0;
    int l = list.length;
    boolean res = true;
    while (index < (l - 1) && res) {
        if (list[index] != list[(l - index) - 1]) {
            res = false;
        }
        index++;
    }
    return res;
}
\end{lstlisting}

El metodo recorre el arreglo desde el extremo izquierdo y compara cada caracter con su
\emph{espejo} en el extremo derecho hasta cruzar el centro o detectar un par desigual.

\subsection{Suite JUnit propuesta}

\begin{lstlisting}[style=codestyle,language=Java]
@Test public void testCapicua() {
    char[] a = {'n','e','u','q','u','e','n'};
    assertTrue(Palindrome.capicua(a));
}
@Test public void testNoCapicuaPar() {
    assertFalse(Palindrome.capicua(new char[]{'a','b'}));
}
@Test public void testNoCapicuaConExtremosIguales() {
    assertFalse(Palindrome.capicua(new char[]{'a','b','c','a'}));
}
@Test public void testArregloVacioEsCapicua() {
    assertTrue(Palindrome.capicua(new char[]{}));
}
@Test public void testUnElementoEsCapicua() {
    assertTrue(Palindrome.capicua(new char[]{'x'}));
}
\end{lstlisting}

\textbf{Por que estos cinco tests.}

\begin{itemize}
  \item \texttt{testCapicua}: caso positivo tipico (capicua impar de varios caracteres).
  \item \texttt{testNoCapicuaPar}: minimo no capicua de longitud par.
  \item \texttt{testNoCapicuaConExtremosIguales}: clave para evitar que la implementacion solo
  compare el primer y ultimo caracter. Si un mutante se queda con la primera comparacion,
  responderia \texttt{true} y este test lo detecta.
  \item \texttt{testArregloVacioEsCapicua} y \texttt{testUnElementoEsCapicua}: cubren los dos
  bordes ($l = 0$ y $l = 1$) que naturalmente devuelven \texttt{true} sin entrar al bucle.
\end{itemize}

\subsection{Reporte de PIT}

\begin{center}
\begin{tabular}{lc}
\toprule
Metrica & Valor \\
\midrule
Mutantes generados & 9 \\
Mutantes muertos   & 8 \\
Mutantes sobrevivientes & 1 \\
\emph{Mutation Score} bruto & $8/9 = 89\%$ \\
\emph{Mutation Score} sobre no equivalentes & $8/8 = 100\%$ \\
\bottomrule
\end{tabular}
\end{center}

Antes de completar la suite el score base (template) era $5/9 = 56\%$. Cada test que se sumo
elimino mutantes especificos asociados al borde del bucle, a la negacion de la condicion del
\texttt{if} y al manejo de arreglos cortos.

\subsection{Respuestas}

\paragraph{(a) \textquestiondown{}Cuantos mutantes genero la herramienta?}
\textbf{9 mutantes} sobre el metodo \texttt{capicua}.

\paragraph{(b) \textquestiondown{}Cuantos tests se necesitaron?}
\textbf{5 tests}, suficientes para matar todos los mutantes \emph{no equivalentes}.

\paragraph{(c) \textquestiondown{}Que puntaje se obtuvo antes de analizar equivalentes?}
\textbf{89\% ($8/9$)}. El unico mutante restante es el que se justifica como equivalente.

\paragraph{(d) Mutantes equivalentes.}

Hay \textbf{1 mutante equivalente}, asociado a la guarda del \texttt{while}:

\begin{lstlisting}[style=codestyle]
Original: while (index <  (l - 1) && res)
Mutado:   while (index <= (l - 1) && res)
\end{lstlisting}

Se analiza por casos sobre la longitud $l$:

\begin{enumerate}
  \item $l = 0$: ambos guardas son falsas ($0 < -1$ y $0 \le -1$). No se entra al ciclo. Ambas
  versiones retornan \texttt{true}.
  \item $l = 1$: original no entra ($0 < 0$ es falso). Mutado entra una vez ($0 \le 0$), compara
  \texttt{list[0]} con \texttt{list[(1-0)-1] = list[0]}, que es la misma posicion: la igualdad
  es trivial y \texttt{res} no cambia. Termina con \texttt{res = true} igual que el original.
  \item $l \ge 2$: el bucle original ejecuta para $\textit{index} = 0, 1, \dots, l-2$. El mutante
  ademas ejecuta una iteracion extra con $\textit{index} = l - 1$, donde se compara
  \texttt{list[l-1]} con \texttt{list[(l - (l-1)) - 1] = list[0]}. Pero ese mismo par
  ($\textit{list[0]}$, $\textit{list[l-1]}$) ya se comparo en la primera iteracion del bucle
  original. La comparacion adicional es por definicion \emph{redundante}: si los extremos eran
  iguales, lo siguen siendo; si eran distintos, \texttt{res} ya valdria \texttt{false} y la
  segunda parte del \texttt{\&\&} corta el ciclo.
\end{enumerate}

Por lo tanto el mutante \emph{no es distinguible} por ningun input y no debe contarse contra el
score: tras descontarlo, el resultado efectivo es $8/8 = 100\%$.

\section{Ejercicio 2: \texttt{TriTyp.triang}}

\subsection{Codigo bajo prueba}

El metodo \texttt{triang(s1,s2,s3)} clasifica un triangulo: 1 escaleno, 2 isosceles, 3
equilatero, 4 no triangulo (lados no positivos o desigualdad triangular incumplida). Su flujo
encadena varios \texttt{if} sobre igualdades de lados, sumas y comparaciones con \texttt{triOut}.

\subsection{Configuracion necesaria de PIT}

El \texttt{pom.xml} apunta por defecto a \texttt{Palindrome}. Para este ejercicio se le indica
explicitamente a PIT cual es la clase y la suite objetivo:

\begin{lstlisting}[style=codestyle]
JAVA_HOME=$(/usr/libexec/java_home -v 17) mvn -Dmaven.repo.local=.m2 \
    -Djacoco.skip=true \
    -DtargetClasses=assignment6_exercises.TriTyp \
    -DtargetTests=assignment6_exercises.TriTypTest \
    test org.pitest:pitest-maven:mutationCoverage
\end{lstlisting}

\subsection{Suite construida}

Para llegar a \emph{Mutation Score} completo hubo que pensar la suite por familias semanticas
de casos. Se proponen \textbf{15 tests}, agrupados asi:

\begin{itemize}
  \item \textbf{Equilatero valido:} \texttt{triang(3,3,3) = 3}.
  \item \textbf{Escaleno valido:} \texttt{triang(3,4,5) = 1}.
  \item \textbf{No triangulo (escaleno) por desigualdad triangular,} en las tres permutaciones
  donde la suma de dos lados iguala o queda por debajo del tercero:
  \texttt{triang(1,2,3) = 4}, \texttt{triang(5,2,3) = 4}, \texttt{triang(2,5,3) = 4}.
  \item \textbf{No triangulo (escaleno) por desigualdad estricta:}
  \texttt{triang(2,3,10) = 4}.
  \item \textbf{Isosceles valido en sus tres variantes:} \texttt{triang(2,2,3) = 2},
  \texttt{triang(2,3,2) = 2}, \texttt{triang(3,2,2) = 2}.
  \item \textbf{Isosceles que no es triangulo} (par igual pero el tercero rompe la
  desigualdad): \texttt{triang(1,1,2) = 4}, \texttt{triang(1,2,1) = 4},
  \texttt{triang(2,1,1) = 4}.
  \item \textbf{Lado no positivo} en cada una de las tres posiciones:
  \texttt{triang(0,2,2) = 4}, \texttt{triang(2,0,2) = 4}, \texttt{triang(2,2,0) = 4}.
\end{itemize}

Cada familia ejercita una decision distinta del metodo (igualdades, sumas, \texttt{triOut} en
$\{0,1,2,3,>3\}$, validez de lados), garantizando que cada operador de mutacion encuentre algun
caso que lo detecte.

\subsection{Reporte de PIT}

\begin{center}
\begin{tabular}{lc}
\toprule
Metrica & Valor \\
\midrule
Mutantes generados & 39 \\
Mutantes muertos   & 39 \\
Mutantes sobrevivientes & 0 \\
\emph{Mutation Score} & $39/39 = 100\%$ \\
\bottomrule
\end{tabular}
\end{center}

\textbf{Progreso:} partiendo de un solo test (\texttt{triang(3,3,3)}) el score base era
$9/39 \approx 23\%$. La adicion de los 14 casos restantes lo llevo al $100\%$.

\subsection{Respuestas}

\paragraph{(a) \textquestiondown{}Cuantos mutantes?} \textbf{39 mutantes.}

\paragraph{(b) \textquestiondown{}Cuantos tests?} \textbf{15 tests} en \texttt{TriTypTest}.

\paragraph{(c) \textquestiondown{}Score antes de analizar equivalentes?} \textbf{$100\%$
($39/39$).}

\paragraph{(d) Mutantes equivalentes.}

\textbf{0 mutantes equivalentes observados.}

La justificacion es directa: la suite mato todos los mutantes generados, por lo que no quedaron
sobrevivientes a clasificar como equivalentes. Esto no significa que el operador de mutacion
sea incapaz de producir equivalentes para este metodo en general; simplemente que, en esta
corrida, ninguno aparecio como sobreviviente. Si hubiera quedado alguno vivo, habria que
analizar caso por caso si admite o no input distinguidor.

\section{Ejercicio 3: Fuzzing}

\subsection{Estructura del paquete \texttt{fuzzing}}

El paquete provee:

\begin{itemize}
  \item \texttt{Fuzzer}: interfaz con un unico metodo \texttt{String fuzz()}.
  \item \texttt{Mutator}: clase utilitaria con operadores de mutacion sobre strings.
  \item \texttt{MutationFuzzer}: fuzzer basado en mutaciones a partir de un conjunto de
  \emph{semillas} validas.
  \item \texttt{RandomFuzzer}: fuzzer puramente aleatorio.
  \item \texttt{HttpProgramRunner}: punto de comparacion --- mide cuantos strings generados por
  cada fuzzer son URLs validas.
  \item \texttt{LinuxCommandTest}: test parametrizado que ejecuta \texttt{bc} sobre $100$
  entradas mutadas y chequea ausencia de crashes.
\end{itemize}

\subsection{(a) \texttt{Mutator}}

Tres operadores de mutacion sobre cadenas, todos protegidos contra entradas \texttt{null} y
manejando explicitamente la cadena vacia donde corresponde:

\begin{enumerate}
  \item \texttt{deleteRandomCharacter(s)} elige una posicion aleatoria
  $\textit{pos} \in [0, |s|)$ y devuelve $s$ sin ese caracter. Si $|s| = 0$, devuelve $s$
  inalterada.
  \item \texttt{insertRandomCharacter(s)} elige $\textit{pos} \in [0, |s|]$ (permite insertar
  al final) y un \texttt{char} aleatorio en el rango ASCII imprimible $[32, 126]$, y lo
  inserta.
  \item \texttt{flipRandomCharacter(s)} elige $\textit{pos} \in [0, |s|)$ y aplica XOR contra
  un solo bit pseudoaleatorio de los siete bits bajos del caracter. Esto produce un caracter
  distinto en la misma posicion. Tambien protege la cadena vacia.
\end{enumerate}

El metodo publico \texttt{mutate(s)} elige uniformemente entre los tres operadores y los
delega.

\textbf{Por que estos tres operadores.} Forman una base minima que cubre los tres modos de
cambio sintactico sobre un string: \emph{eliminar} un simbolo, \emph{agregar} uno, y
\emph{modificar} uno existente. Permiten alcanzar progresivamente toda cadena finita en pocos
pasos sucesivos (encadenando \texttt{mutate} sobre el resultado anterior).

\subsection{(b) \texttt{MutationFuzzer}}

\begin{lstlisting}[style=codestyle,language=Java]
public String fuzz() {
    String candidate = population.get(random.nextInt(population.size()));
    int mutations = min_mutations;
    if (max_mutations > min_mutations) {
        mutations += random.nextInt((max_mutations - min_mutations) + 1);
    }
    for (int i = 0; i < mutations; i++) {
        candidate = Mutator.mutate(candidate);
    }
    return candidate;
}
\end{lstlisting}

\textbf{Funcionamiento.}

\begin{enumerate}
  \item Elige una semilla al azar de la poblacion (cadenas inicialmente validas).
  \item Determina cuantas mutaciones aplicar en este \texttt{fuzz()}, entre
  \texttt{min\_mutations} y \texttt{max\_mutations} (inclusivo).
  \item Aplica las mutaciones \emph{en cadena}: cada paso opera sobre el resultado anterior.
\end{enumerate}

\textbf{Correccion vs. el template original.} El constructor del template tenia un bug por el
cual \texttt{this.max\_mutations} no se asignaba correctamente; el codigo final asigna los tres
parametros explicitamente. Tambien se agregaron precondiciones (\texttt{population} no vacia,
rango de mutaciones consistente).

\subsection{(c) \texttt{RandomFuzzer}}

\begin{lstlisting}[style=codestyle,language=Java]
public String fuzz() {
    int stringLength = random.nextInt(maxLength + 1);
    StringBuilder result = new StringBuilder(stringLength);
    for (int i = 0; i < stringLength; i++) {
        char c = (char) (charStart + random.nextInt(charRange));
        result.append(c);
    }
    return result.toString();
}
\end{lstlisting}

\textbf{Funcionamiento.}

\begin{enumerate}
  \item Elige una longitud aleatoria $L \in [0, \texttt{maxLength}]$.
  \item Construye una cadena de longitud $L$ donde cada caracter es aleatorio en
  $[\texttt{charStart}, \texttt{charStart} + \texttt{charRange})$.
\end{enumerate}

Es el contrapunto del \texttt{MutationFuzzer}: no parte de ninguna entrada valida y por lo
tanto rara vez produce strings que ``parezcan'' validos.

\subsection{(d) \texttt{LinuxCommandTest}}

El test parametrizado ejecuta \texttt{bc} (calculadora Linux) sobre $100$ cadenas generadas por
el fuzzer de mutacion (que parte de la semilla \texttt{"2 + 2"}).

\begin{lstlisting}[style=codestyle,language=Java]
@ParameterizedTest
@MethodSource("stringProvider")
public void bcCommandTest(String data) throws IOException, InterruptedException {
    String interactiveData = data + "\nquit";
    writer.write(interactiveData); writer.close();

    ProcessBuilder pb = new ProcessBuilder(program, FILE);
    Process process = pb.start();

    // ... captura stdout, stderr ...

    int exitCode = process.waitFor();

    assertNotEquals(134, exitCode); // abort
    assertNotEquals(139, exitCode); // segfault
    String stderrLower = stderr.toString().toLowerCase();
    assertFalse(stderrLower.contains("segmentation fault"));
    assertFalse(stderrLower.contains("core dumped"));
    assertFalse(stderrLower.contains("illegal instruction"));
}
\end{lstlisting}

\textbf{Diseno del oraculo.}

\begin{itemize}
  \item \textbf{No se exige \texttt{exitCode == 0}.} Es normal que \texttt{bc} rechace entradas
  malformadas con error sintactico; eso \emph{no es un crash}.
  \item \textbf{Si se rechazan signales tipicos de fallo grave:} exit code $134$ (SIGABRT),
  $139$ (SIGSEGV), o mensajes de \texttt{stderr} que mencionen segmentation fault, core dumped
  o illegal instruction.
\end{itemize}

La consigna prohibe expresamente cambiar \texttt{bc} por otro programa, ya que el fuzzer puede
generar entradas peligrosas si se las pasara a un programa que modifica el sistema. \texttt{bc}
es seguro: solo lee expresiones aritmeticas.

\subsection{Ejecucion}

\begin{lstlisting}[style=codestyle]
JAVA_HOME=$(/usr/libexec/java_home -v 17) mvn -Dmaven.repo.local=.m2 \
    -Djacoco.skip=true test
\end{lstlisting}

\textbf{Resultado:} \texttt{BUILD SUCCESS}. \texttt{LinuxCommandTest} pasa los $100$ casos
parametrizados y \texttt{TriTypTest} pasa sus $15$ tests.

\section*{Conclusiones}
\addcontentsline{toc}{section}{Conclusiones}

\begin{itemize}
  \item \textbf{Testing de mutacion} es una herramienta poderosa para evaluar la
  \emph{capacidad de deteccion} de una suite: no alcanza con tener cobertura de lineas o de
  ramas. PIT inyecta variaciones sintacticas tipicas (cambiar operadores, eliminar
  asignaciones, mover constantes, negar condicionales), y un test debe ser capaz de
  distinguirlas. En el caso de \texttt{Palindrome}, una suite de 5 tests bien pensados llega a
  $89\%$ y deja un unico mutante \emph{equivalente} identificable analiticamente.
  \item En \texttt{TriTyp} el score alcanzo el $100\%$, pero esto fue resultado de una suite de
  15 tests \emph{enfocada} en las decisiones del codigo. Una cobertura ingenua (e.g.\ todos
  equilateros) deja la enorme mayoria de los mutantes vivos.
  \item \textbf{Fuzzing} es una tecnica complementaria: en lugar de pedir tests humanos, genera
  miles de entradas. \texttt{MutationFuzzer} y \texttt{RandomFuzzer} muestran el tradeoff entre
  \emph{validez} y \emph{exploracion}: las mutaciones de semilla preservan estructura (mas
  inputs validos), la generacion aleatoria explora mas pero acierta menos. La eleccion del
  oraculo (que considera un ``fallo'') es tan importante como la generacion misma.
\end{itemize}

\end{document}
